% Cover letter - CFC Associate Data Scientist
\documentclass[letterpaper,11pt]{article}
\usepackage[empty]{fullpage}
\usepackage[hidelinks]{hyperref}
\usepackage[english]{babel}
\usepackage{parskip}
\setlength{\parskip}{6pt}
\usepackage{setspace}
\addtolength{\oddsidemargin}{-0.5in}
\addtolength{\evensidemargin}{-0.5in}
\addtolength{\textwidth}{0.95in}
\addtolength{\topmargin}{-.80in}
\addtolength{\textheight}{1.7in}
\urlstyle{same}
\raggedright
\setstretch{1.0}
\pagestyle{empty}
\begin{document}

Dear Hiring Manager,

I am writing to apply for the Associate Data Scientist position in the Data \& AI team at CFC. I am currently completing an MSc in Health Data Science at University College London, where a distinction is expected, after graduating from Imperial College London, and I work as a Data Analyst Intern at TikTok Shop UK. Across my academic and professional experience I have worked on both sides of the problem this role describes: building statistical and machine learning models, and building the evaluation frameworks that establish whether an AI system actually holds up outside a notebook. I am particularly interested in this role because it combines rigorous experimentation with production AI systems that carry real consequences.

What attracts me most to CFC is the opportunity to sit as the analytical counterpart to engineers building agentic systems, running the experiments and generating the evidence that decides what gets built. I enjoy taking an unclear business problem, identifying what users actually need, and turning it into a workf\kern0pt low that can be tested and improved. The description of careful, evidence-based thinking as the dif\kern0pt ference between an AI system that works in a demonstration and one that survives real conditions matches how I have been taught to work and how I prefer to work.

At TikTok Shop UK I explore large and messy merchant and operations datasets in Python and SQL, and apply statistical analysis to explain what is driving performance across the UK marketplace. I build KPI frameworks and dashboards that operations teams use to set targets, and I present f\kern0pt indings and recommendations to product, operations and commercial stakeholders who do not want the technical detail. I also built an AI agent, using Claude Code, that automates repetitive steps of my own analytics workf\kern0pt low, and my team and adjacent teams now use it.

At DS Digital Technology I worked on a medical AI engagement for a pharmaceutical client, and that work is the closest match to this role. I designed a 246-question benchmark and a six-dimension safety rubric with red-line rules, agreed with subject-matter experts, which became the client's acceptance standard for evaluating patient-facing AI answers. I then built a Python pipeline calling six LLM APIs with retrieval and search augmentation, which produced and scored more than 2,400 responses through an LLM-as-judge framework. I co-built a medically reviewed knowledge graph of 123 entities and 161 relations and designed blinded paired A/B tests to measure whether it improved answer quality, rather than assuming that it did. Much of the value of that project was in f\kern0pt inding where the systems failed and why.

My MSc dissertation is where the statistical side was built. I compared logistic regression against three tree ensembles for cardiovascular risk prediction on 458,840 UK Biobank participants, with attention to calibration rather than headline accuracy, then used doubly robust causal forests on a longitudinal subset to estimate the ef\kern0pt fect of lifestyle change. I tested how far the conclusions could be pushed before they broke, using overlap weighting, sensitivity analysis including E-values, and external validation on a second cohort, and I reported null results where I found them. I delivered the work as a Streamlit application so that people could run scenarios without touching the code.

My time at Ipsos Strategy3 taught me the communication side. I structured and analysed data from more than 80 expert interviews, built market-sizing and forecasting models, and turned the results into dashboards and recommendations for client stakeholders. This experience has strengthened my ability to explain an analytical decision to an audience that needs the conclusion rather than the method.

I should be clear about where I am starting from. My production engineering experience is limited: I have not yet owned a system of my own running in production, and my AWS experience is at a basic level, so the productionisation and monitoring side of the role is where I would need support early on. My Python, my SQL, my statistics and my experience designing and running evaluations of LLM-based systems are well established. I would bring CFC an analyst who is comfortable with rigorous experimental design and with generative AI evaluation, and who is used to working in a small multidisciplinary team. I have the right to work in the UK until January 2029 without sponsorship, and I am based in London.

Thank you for considering my application. I would welcome the opportunity to discuss how my experience in LLM evaluation, experimentation and statistical analysis could support the Data \& AI team.

\vspace{2pt}
Yours sincerely,

\vspace{8pt}
Zhuoyan Sun

\end{document}
